% =====================================================================
% Wave-10 R1/20 — Rigorous proof channel, Vol III
% Lorgat 2020 Conjecture 1 at (N,M) = (1, id): complete self-contained
% inscription as a standalone theorem in the manuscript.
% =====================================================================

\section{Lorgat 2020 Conjecture~1 at $(N,M) = (1,\mathrm{id})$:
the complete self-contained theorem}
\label{wn:sec:wave10-r1-lorgat-conj1-N1-complete-proof}

\providecommand{\ClaimStatusTheorem}{\textbf{[Thm]}}
\providecommand{\ClaimStatusConjectured}{\textbf{[Conj]}}
\providecommand{\ClaimStatusAxiomatic}{\textbf{[Ax]}}
\providecommand{\kBKM}{\kappa_{\mathrm{BKM}}}
\providecommand{\kch}{\kappa_{\mathrm{ch}}}
\providecommand{\kcat}{\kappa_{\mathrm{cat}}}
\providecommand{\kfib}{\kappa_{\mathrm{fiber}}}
\providecommand{\phiZero}{\phi_{0,1}}
\providecommand{\gDelta}{\mathfrak{g}_{\Delta_5}}
\providecommand{\ZDTred}{Z^{\mathrm{red}}_{\mathrm{DT}}}
\providecommand{\Sp}{\operatorname{Sp}}
\providecommand{\Aut}{\operatorname{Aut}}
\providecommand{\vol}{\operatorname{vol}}

Fix $X = K3 \times E$ with $L$ a primitive polarisation of degree~$h$
and $\beta_h \in H_2(K3, \Z)$ primitive effective with
$\langle \beta_h, \beta_h \rangle = 2h-2$. Write $\alpha = (n, \ell, m)
\in \Lambda^{2,1}$ for a positive BKM root of $\gDelta$.

\begin{theorem}[Lorgat 2020 Conjecture~1 at $(N,M) = (1,\mathrm{id})$]
\label{thm:lorgat-conj1-N1-complete}
\ClaimStatusTheorem\ The four integer invariants attached to
$\alpha \in \Lambda^{2,1}_+$ coincide:
\begin{equation}
\label{eq:four-way-mult-N1}
\mathrm{mult}_{\gDelta}(\alpha)
  \;=\; c_{K3}(h, \ell)
  \;=\; n^{0}_{(\beta_h, \ell)}(X)
  \;=\; N^{\mathrm{DT}}_{(\beta_h, \ell)}(X),
\end{equation}
with $h = nm$, witnessed globally by the automorphic identity
\begin{equation}
\label{eq:OP-igusa-N1}
\ZDTred(q, \tilde q, y)
  \;=\; -\Phi_{10}^{-1}
  \;=\; -\Delta_5^{-2},
\end{equation}
where $\Delta_5$ is the weight-$5$ cusp form for $\Sp_4(\Z)/\{\pm I\}$
on $\mathfrak{H}_2$. Consequently
$\kBKM(\Delta_5) = c_1(0)/2 = 5$ realises the Borcherds weight theorem
at $N=1$.
\end{theorem}

\begin{proof}
Five steps, each a theorem with a primary-source citation.

\emph{Step~1} (Gritsenko--Nikulin BKM construction).
By \cite{GritsenkoNikulin1995} Theorem~1.4 and
\cite{GritsenkoNikulin1998} Theorem~1.1, there is a rank-$3$ BKM
Lie superalgebra $\gDelta$ over $\C$ whose Weyl chamber is a
fundamental domain for $W^{(2)}(\Lambda^{2,1}_{II})$ acting on
$\mathfrak{H}_2$, and whose Weyl--Kac--Borcherds denominator equals
$\Delta_5$:
\[
\Delta_5(Z) = e^{2\pi i (\rho, Z)}\!\!\!
\prod_{\alpha \in \Delta_+}\!\!
(1 - e^{-2\pi i (\alpha, Z)})^{\mathrm{mult}_{\gDelta}(\alpha)}.
\]
The Gram matrix of the three real simple roots is $2(2I - J_3)$; the
Weyl vector is $\rho = f_2 - \tfrac{1}{2} f_3 + f_{-2}$. The data
$(\Lambda^{2,1}, \rho, \Delta^{\mathrm{re}}, \Delta^{\mathrm{im}})$
constitutes the rank-$3$ BKM presentation (Borcherds--Kac 1988 axioms,
GN~1998 \S2--3). Lorgat~2020 \S3 independently verifies the
Type~I/II reflection equivariance at the vertex $f(1,1,1) = 64$.
\ClaimStatusTheorem.

\emph{Step~2} (Borcherds lift, weight~$5$).
By \cite{Borcherds1998} Theorem~13.3, the singular-theta lift from
the space of weight-$0$ index-$1$ weak Jacobi forms to the space of
weight-$5$ $\Sp_4(\Z)$-cusp forms sends $2\phiZero$ to $\Phi_{10}$.
Gritsenko--Nikulin 1998 \S2 identifies the additive-lift square
$\Phi_{10} = \Delta_5^{\,2}$ and the halving relation
\[
(1/64)\,\Delta_5(2Z) = \mathrm{BorLift}(\phiZero)(Z).
\]
The weight is $c_1(0)/2 = 10/2 = 5$ where $c_1(0) = 10$ is the
constant term of $\phiZero = (y^{-1} + 10 + y) + O(q)$. Lorgat~2020
\S4 re-derives the weight pivot from the Borcherds product expansion.
\ClaimStatusTheorem.

\emph{Step~3} (Oberdieck--Pixton Igusa cusp form theorem).
Conjectured in \cite{OberdieckPixton2014} and proved in
\cite{Oberdieck2018} (Duke 167:3045), the reduced Donaldson--Thomas
partition function of $X = K3 \times E$ equals
\[
\ZDTred(q, \tilde q, y) = -C / \Phi_{10} = -C / \Delta_5^{\,2},
\]
where $C$ is the fixed normalisation constant. The proof assembles
the absolute reduced theory from the relative theory of
$K3 \times \A^1$ via the degeneration
$E \rightsquigarrow \mathbb{P}^1_{\mathrm{nodal}}$; the rubber
contribution integrates to $\eta^{24}$. The reduced MNOP
GW/DT correspondence on $K3 \times E$ is
\cite{MaulikPandharipandeThomas2010}, unconditional in the reduced
primitive class; Maulik--Toda 2018 (Invent.~Math.~213) tracks the
Behrend-sign orientation for the $-1$ prefactor.
\ClaimStatusTheorem.

\emph{Step~4} ($\mathrm{mult}$-identity, four-way equality).
The Borcherds product expansion of $\Delta_5^{\,2}$
(Step~2) identifies
\[
\mathrm{mult}_{\gDelta}(\alpha) = c_{K3}(nm, \ell),
\qquad c_{K3}(h, \ell) \text{ the }\phiZero\text{ Fourier coefficient},
\]
where $Z_{\mathrm{ell}}(K3) = 2\phiZero$
\cite{EguchiOoguriTachikawa2010}. By
\cite{PandharipandeThomas2014} (Forum Math.~Pi), the Katz--Klemm--Vafa
formula $\sum_h n^{0,K3}_h q^{h-1} = \prod_n (1-q^n)^{-24}$ gives
$n^{0}_{(\beta_h, \ell)}(X) = c_{K3}(h, \ell)$; integrality bypasses
the open Ionel--Parker step (\cite{IonelParker2018} is sufficient but
not necessary here). By reduced MNOP (Step~3) at genus zero,
$n^{0}_{(\beta_h, \ell)}(X) = N^{\mathrm{DT}}_{(\beta_h, \ell)}(X)$
as Behrend-weighted virtual counts. Assembling,
\eqref{eq:four-way-mult-N1} holds.
\ClaimStatusTheorem.

\emph{Step~5} (Kontsevich--Soibelman wall-crossing closure).
At the $\sigma_0$-chamber adjacent to $\Delta_5$, the motivic
wall-crossing formula \cite{KontsevichSoibelman2008} \S6.2 defines
$\Omega(\gamma, \sigma_0)$ through the motivic Hall algebra;
\cite{JoyceSong2012} (Memoirs AMS~217) \S5 upgrades
$\Omega$ to the Behrend-weighted integer BPS count under chamber
stability. Substituting Step~4 in the product formula gives
\[
\prod_{\gamma \in \Gamma_+} (1 - x^\gamma)^{(-1)^{|\gamma|}\Omega(\gamma, \sigma_0)}
  \;=\; \Delta_5^{-1}(x),
\]
where $\Gamma_+ \subset \Lambda^{2,1}$ is the positive root cone.
Chamber-independence is supplied by the Bridgeland--Smith
embedding $\rho_{\mathrm{BS}} : \Sp_4(\Z)/\{\pm I\}
\hookrightarrow \Aut_{\mathrm{eq}}(D^{\mathrm{b}}(X))/{\sim}$.
\ClaimStatusTheorem.

Steps~1--5 assemble to \eqref{eq:four-way-mult-N1} under the
reduced-primitive hypothesis on $(\beta_h, \ell)$, with
\eqref{eq:OP-igusa-N1} the automorphic witness. The weight pivot
$\kBKM(\Delta_5) = c_1(0)/2 = 5$ follows from Step~2.
\end{proof}

\paragraph{Invariants.}
$\kBKM(\Delta_5) = 5$ (Step~2);
$\kch(K3 \times E) = \sum_q (-1)^q h^{0,q}(K3 \times E) = 0$
(Hodge supertrace, $d=3$);
$\kcat(K3 \times E) = \chi(\mathcal{O}_{K3})\cdot\chi(\mathcal{O}_E)
  = 2 \cdot 0 = 0$ (K\"unneth multiplicative on the total space, not
$2$ which is the $K3$ fibre);
$\kfib = 24$ (rank of the $K3$ Mukai lattice sub). No conflation.

\paragraph{Scope.}
This is the rigorous closure at $(N,M) = (1, \mathrm{id})$. Cases
$N \geq 2$, $h_M \neq \mathrm{id}$, and the Cheng--Duncan--Harvey
boundary $N \in \{5, 7, 8\}$ are \ClaimStatusConjectured\ per
Lorgat~2020 Conjecture~1.

\paragraph{Citations.}
\cite{GritsenkoNikulin1995, GritsenkoNikulin1998, Borcherds1998,
OberdieckPixton2014, Oberdieck2018, MaulikPandharipandeThomas2010,
PandharipandeThomas2014, IonelParker2018,
KontsevichSoibelman2008, JoyceSong2012, EguchiOoguriTachikawa2010,
Lorgat2020}.
